\section{Security Analysis}

\label{sec:security}

\subsection{Threat Model}

We consider the following adversaries:

\begin{enumerate}[leftmargin=1.4em]
    \item \textbf{Lazy worker}: Skips computation, returns random or cached results.
    \item \textbf{Colluding workers}: Multiple workers coordinate to produce consistent but incorrect results.
    \item \textbf{Sybil attacker}: Creates many fake worker identities to dominate the network.
    \item \textbf{TEE attacker}: Compromises TEE hardware to forge attestations.
    \item \textbf{Byzantine validator}: Votes to accept incorrect results.
\end{enumerate}

\subsection{Formal Security Proofs}

\begin{theorem}[Lazy Worker Detection]
A lazy worker who skips computation on fraction $\delta$ of jobs is detected within $T$ epochs with probability at least $1 - (1-\alpha)^{\delta N T}$, where $\alpha$ is the sampling rate and $N$ is jobs per epoch.
\end{theorem}

\begin{proof}
Each job is independently sampled with probability $\alpha$. The worker produces incorrect output on $\delta N$ jobs per epoch. Each incorrect job is caught with probability $\alpha$. Over $T$ epochs, the probability of catching at least one is $1 - (1-\alpha)^{\delta N T}$. For $\alpha = 0.003$, $\delta = 0.01$, $N = 10{,}000$, $T = 10$: $P = 1 - (0.997)^{1000} = 0.9503$.
\end{proof}

\begin{theorem}[Collusion Resistance]
For $k$ colluding workers, the expected cost of successful collusion exceeds the expected gain when:
\begin{equation}
k \cdot \text{Stake} \cdot P(\text{detect}) \cdot \text{SlashRate} > k \cdot \text{Savings}_{\text{lazy}}.
\end{equation}
\end{theorem}

\begin{proof}
Colluding workers cannot reduce their collective detection probability below $1 - (1-\alpha)^{\delta N T}$ because verification is performed by independent verifier nodes. The expected loss from slashing ($\text{Stake} \times P(\text{detect}) \times \text{SlashRate}$) exceeds the savings from lazy computation when the stake requirement satisfies:
$\text{Stake} > \text{Savings}_{\text{lazy}} / (P(\text{detect}) \times \text{SlashRate})$.
For typical parameters, this requires $\text{Stake} > 4.2 \times \text{Revenue}_{\text{epoch}}$.
\end{proof}

\begin{theorem}[Sybil Resistance]
The cost of a Sybil attack that controls fraction $f$ of the network's compute capacity is at least $f \cdot N_{\text{workers}} \cdot \text{MinStake}$, where $\text{MinStake}$ is the minimum stake per worker.
\end{theorem}

\begin{proof}
Each Sybil identity must stake $\text{MinStake}$ tokens. To control fraction $f$ of capacity, the attacker needs $f \cdot N_{\text{workers}}$ identities (assuming uniform capacity per worker). The total stake cost is linear in $f$.
\end{proof}

\subsection{Economic Incentive Analysis}

\begin{proposition}[Nash Equilibrium]
Honest computation is a Nash equilibrium of the PoAI game when:
\begin{equation}
\text{Revenue}_{\text{honest}} > \text{Revenue}_{\text{lazy}} + \text{Expected\_Slash},
\end{equation}
which holds when $P(\text{detect}) \cdot \text{Slash} > \text{Savings}_{\text{lazy}} \cdot (1 - P(\text{detect}))$.
\end{proposition}

For ACI's parameters ($\alpha = 0.003$, 100\% slash on detection, stake = $5\times$ monthly revenue):

\begin{equation}
P(\text{detect}) \cdot 5R > \delta R \cdot (1 - P(\text{detect})),
\end{equation}

which is satisfied for any $\delta > 0$ over sufficiently many epochs, making honest computation the dominant strategy.

